\section{Related Work}\label{sec:related}

\paragraph{Specification mining from traces.}
Flie~\cite{NeiderG18} learns the smallest LTL formula separating
positive from negative lasso traces via a SAT encoding over syntax
DAGs---our propositional encoding builds on this framework but
eliminates temporal operator search by fixing the \GRone skeleton.
In the finite-trace setting,
Scarlet~\cite{RahaRFV22} learns LTL$_f$ formulas via scalable
anytime algorithms, and
Bolt~\cite{BathieFMMV26} reduces LTL$_f$ learning to Boolean set
cover.
Fijalkow and Lagarde~\cite{FijalkowL21} establish complexity bounds
for passive LTL learning, showing NP-hardness even for restricted
fragments.
All of these approaches mine monolithic temporal formulas without
distinguishing environment assumptions from system guarantees,
making the resulting specifications unsuitable for reactive synthesis
without manual restructuring into an assume--guarantee form.

\paragraph{Constrained specification mining.}
Texada~\cite{LemieuxDB15} was among the first to introduce
user-specified property templates into specification mining, allowing
users to constrain the syntactic shape of mined properties.
LTL-Sketcher~\cite{LutzNR23} extends this idea to sketch
completion: given a partial LTL formula with holes, it fills them
from positive and negative examples using a SAT encoding.
ATLAS~\cite{ZhangCGD25} further generalises structural constraints
via first-order relational logic in Alloy, reducing constrained LTL
learning to MaxSAT.
One can encode \GRone structure as an ATLAS constraint
(Appendix~\ref{app:alloy}), but the underlying encoding retains
full LTL temporal semantics, treating \GRone as a filter on the
general search space rather than exploiting it algorithmically.
Orthogonally, work on \GRone assumption
mining~\cite{ChatterjeeHJ08,LiDS11} automatically generates or
refines environment assumptions for a \emph{given} set of system
guarantees, while specification
debugging~\cite{KonighoferHB13} localises errors in existing
\GRone specifications.
These approaches assume that part of the specification is already
known; \toolname mines the complete assume--guarantee formula from
traces alone.
More broadly, \toolname differs from all of the above in that the
\GRone template is not a constraint layered atop a general LTL miner
but the foundation of the encoding itself: temporal semantics are
hardcoded, the search is purely propositional, and incremental
solving shares work across template configurations.
